\section{Controlled densities and a common continuation moment set}\label{sec:controlled51}
The reset construction has an action-independent transition law. We now allow the action to change the entire next-state density. The example also permits continuous dependence on the current state and optimizes over the original Borel policy class. A finite-dimensional common continuation, rather than a discretization of that class, makes the problem tractable.

Consider two decision dates on $X=[0,1]$, a uniform initial law, zero terminal reward, and actions $a=0,1$. Action zero is installed. Let $h,g,k$ be strictly positive affine functions, $h\ge\varepsilon$, $c>0$, and $|\theta_a|\le1$. Write $\phi(x)=2x-1$ and specify
\begin{equation}\label{eq:kernel51}
 P_0^a(dy\mid x)=[1+\theta_a\phi(x)\phi(y)]\,dy.
\end{equation}
This is a nonnegative, normalized continuous density, not a controlled atomic kernel. At date one, operating rewards are $(0,h(y))$ and implementation costs are $(0,ch(y))$. At date zero, rewards are $(-\beta P_0^0h-g,-\beta P_0^1h)$ and implementation costs are $(0,k(x))$. Hence $V_1=h$, $V_0=0$, and operating disadvantages are $(h,0)$ at date one and $(g,0)$ at date zero. These are designed economic primitives. Their construction supplies an exactly solvable operating oracle, not evidence about the cost of a learned oracle.

For any common continuation let $D(y)=D_1^p(y)$ and define
\[
 u=\varepsilon^{-1}\int_0^1D(y)\,dy,\qquad
 v=\varepsilon^{-1}\int_0^1\phi(y)D(y)\,dy.
\]
The implementation identity $C_1^p(y)=c[h(y)-D(y)]$ is substantive: it follows from the proportional date-one cost, and it is not imposed on an arbitrary economic model. Set $A=\int h=h(0)+h'/2$ and $B=\int\phi h=h'/6$.

\begin{theorem}[Exact common-moment reduction]\label{thm:controlled51}
The feasible date-one moment set for all Borel policies is
\begin{equation}\label{eq:lens51}
 \mathcal Z=\{(u,v):0\le u\le1,\ |v|\le u(1-u)\}.
\end{equation}
For $z=(u,v)\in\mathcal Z$, define the affine functions
\begin{align}
 e_a(x;z)&=g(x)\ind\{a=0\}+\beta\varepsilon[u+\theta_a\phi(x)v],\label{eq:row51a}\\
 q_a(x;z)&=k(x)\ind\{a=1\}+\beta c[A-\varepsilon u+\theta_a\phi(x)(B-\varepsilon v)].\label{eq:row51b}
\end{align}
Then the original common-policy optimum satisfies
\begin{equation}\label{eq:reduction51}
 \KR=\min_{z\in\mathcal Z}F(z),\qquad
 F(z)=\int_0^1\min_{p\in\Delta_2:\,p\cdot e(x;z)\le\varepsilon}p\cdot q(x;z)\,dx.
\end{equation}
Every feasible moment has an explicit Borel continuation, and every minimizing row has a Borel selection. Thus both directions of this equality concern the original policy object, not a relaxation or an initial lottery over plans.
\end{theorem}
\begin{proof}
Since $h\ge\varepsilon$, every Borel $D$ with $0\le D\le\varepsilon$ is realized by the installed-action probability $D/h$. Write $f=D/\varepsilon$. For $\int f=u$, the maximum of $\int\phi f$ is attained by $f=\ind_{[1-u,1]}$ and equals $u(1-u)$. Indeed, subtracting this indicator and multiplying by $\phi(y)-\phi(1-u)$ gives a nonpositive integrand everywhere; its integral proves the inequality because the two functions have the same mass. The lower bound follows by reflection.

Conversely, when $0<u<1$, put
\[
 \alpha=\frac12\left(1+\frac{v}{u(1-u)}\right),\qquad
 D(y)=\varepsilon[\alpha\ind_{[1-u,1]}(y)+(1-\alpha)\ind_{[0,u]}(y)].
\]
The two indicators may overlap; their coefficients sum to one, so $D\le\varepsilon$ even on the overlap. This function has precisely the stated moments. The cases $u=0,1$ use $D=0,\varepsilon$. Values at interval endpoints can be chosen without changing any moment and without violating the pointwise bound.

The density in \eqref{eq:kernel51} now gives $P_0^aD=\varepsilon[u+\theta_a\phi(x)v]$ and $P_0^aC_1=c[A-\varepsilon u+\theta_a\phi(x)(B-\varepsilon v)]$. These are the same successor moments for both actions and every restart. Conditional on them, the remaining operating restriction is the row constraint in \eqref{eq:reduction51}; implementation expense is its objective. Any feasible original policy therefore costs at least $F(z)$. Conversely, the displayed continuation and a minimizing date-zero row attain that value. Action one always has regret at most $\beta\varepsilon$, leaving uniform slack $(1-\beta)\varepsilon$. A row minimizer is selected from two pure actions and their binding mixture by a fixed tie rule. These are Borel functions of $x$. The slack also gives continuity of the row value in $z$ by mixing any perturbed row toward action one. Uniform boundedness permits integration, and compactness of $\mathcal Z$ yields attainment.
\end{proof}

The moment inequalities are elementary extremal-integral inequalities; the new role they play here is to characterize the \emph{common} operating and implementation continuation simultaneously. Allowing a different $v$ on the incoming edges would invalidate the equality. General finite-rank kernels would require additional common moments and a realizability description; proportional terminal costs and two dates are the executed structure, not conclusions inferred for general horizons.

\subsection{An integrated, two-coordinate global cover}\label{sec:cover51}
For a rectangle $B=[\underline u,\overline u]\times[\underline v,\overline v]$, minimize each $e_a(x;z)$ and $q_a(x;z)$ over its four moment corners independently, obtaining $\underline e_a^B(x)$ and $\underline q_a^B(x)$. The independent minimizations are a relaxation used only for the lower endpoint. Define
\begin{equation}\label{eq:boxlower51}
 b(B)=\int_0^1\min_{p\in\Delta_2:\,p\cdot\underline e^B(x)\le\varepsilon}
 p\cdot\underline q^B(x)\,dx.
\end{equation}
Then $b(B)\le F(z)$ for every $z\in B\cap\mathcal Z$. A rectangle is outside $\mathcal Z$ if its $v$ interval misses $[-M,M]$, where
\[
 M=\max_{u\in[\underline u,\overline u]}u(1-u).
\]
This maximum is rational for rational endpoints: project $1/2$ onto the interval. The test is necessary and sufficient for an empty rectangle--moment-set intersection.

The corner minima switch only at $x=1/2$. On either side, partition further at the rational roots of $\underline e_a^B-\varepsilon$, $\underline e_0^B-\underline e_1^B$, and $\underline q_0^B-\underline q_1^B$. On each open subinterval, the optimal row is a pure action or a mixture with installed probability
\[
 p_0(x)=\frac{\varepsilon-e_1(x)}{e_0(x)-e_1(x)}.
\]
The corresponding integrand is affine or a quadratic polynomial divided by an affine function. Polynomial division leaves a rational term and a rational multiple of a logarithm of a positive rational number. No sample average, floating quadrature, or spatial grid enters a certified endpoint.

To bound logarithms, first reduce the argument by powers of two to $r\in[1,2]$ and let $s=(r-1)/(r+1)$. For every positive integer $N$,
\begin{equation}\label{eq:log51}
 0\le \log r-2\sum_{j=0}^{N-1}\frac{s^{2j+1}}{2j+1}
 \le\frac{2s^{2N+1}}{(2N+1)(1-s^2)}.
\end{equation}
Negative rational multipliers reverse the endpoints. The implementation rounds each series term down and adds its explicit rounding allowance. The constructor uses 26 terms and 80-bit term rounding; the separate reader uses 34 terms and 90-bit term rounding. Final constructor integrals are rounded outward to 48-bit dyadics. At a removable zero denominator the polynomial expression is used directly. Isolated state breakpoints have zero initial mass; the deployed action there is the feasible operating-best action.

\begin{proposition}[Paired cover and quantitative refinement]\label{prop:cover51}
Retain a finite binary cover of $[0,1]\times[-1/4,1/4]$. Discard rectangles only by the exact moment-set test or a lower endpoint exceeding a verified feasible upper cost. If $\ell_B\le b(B)$ is the integrated lower endpoint on every retained rectangle, and $U\ge F(\widehat z)$ for some $\widehat z\in\mathcal Z$, then
\[
 \min\{U,\min_B\ell_B\}\le\KR\le U.
\]
Let $w$ be the sum of a rectangle's side widths and $e_B$ its integral-enclosure width. Put $Q=\max_x k(x)+\beta c\max_y h(y)$. The center followed by projection of $v$ onto $[-u(1-u),u(1-u)]$ is feasible, and its cost obeys
\begin{equation}\label{eq:rate51}
 F(\widehat z_B)-\ell_B
 \le 2\beta\varepsilon c w+
 Q\frac{2\beta w}{1-\beta+2\beta w}+e_B.
\end{equation}
Consequently, fair widest-coordinate bisection and tightening the integration enclosures give an executable stopping rule at every positive target. With fixed positive slack, covering the two moment coordinates requires $O(\eta^{-2})$ rectangles for width $\eta$, apart from arithmetic precision work. This bound concerns this structured two-period model, not general controlled MDPs.
\end{proposition}
\begin{proof}
Every original continuation has one point of $\mathcal Z$. A retained rectangle containing it supplies a valid lower bound by \eqref{eq:boxlower51}; inherited ancestor bounds can only strengthen it. Feasible moments and their constructed Borel policies supply the upper bound. Binary partition identities certify that no feasible moment has been omitted.

For the modulus, $u(1-u)$ is 1-Lipschitz on $[0,1]$. Projecting the center's $v$ coordinate therefore moves it outside the original rectangle by at most its $u$ width. Relative to every corner used in its lower envelope, the sum of the candidate's coordinate distances is at most $2w$. For the relaxed minimizing row, evaluating regrets at the candidate increases them by at most $\delta=2\beta\varepsilon w$, and costs by at most $c\delta$. Mix this row toward action one with weight $\delta/[\delta+(1-\beta)\varepsilon]$. The mixture is feasible and increases its cost by at most that weight times $Q$, since original costs are nonnegative and at most $Q$. The optimal row at the candidate costs no more. Integrate and add $e_B$ to obtain \eqref{eq:rate51}. The right side vanishes with rectangle width and integration error. Uniform bisection in two coordinates yields the stated sufficient count; best-lower-first selection with widest-coordinate splitting processes only a subset of that sufficient finite refinement before its stopping test succeeds.
\end{proof}

A model with $\theta_0=\theta_1=0$ is independent of $v$ and is reduced exactly to one coordinate. An unrestricted run otherwise retains two common moments. For economic interpretation, a separate run fixes $v=0$. Every feasible such pair is realized by constant terminal regret $D(y)=\varepsilon u$, and all policies with that pair have the same continuation value. Thus it exactly solves the counterfactual class with spatially uniform terminal regret. If $[L_0,U_0]$ is its interval and $[L,U]$ the unrestricted interval, $L_0-U>0$ proves a strict implementation saving from allowing state-dependent continuation. It is not a comparison against an external optimization package.
